% !TEX root = ../notes.tex

\documentclass[../notes.tex]{subfiles}

\begin{document}

\section{April 27}
Today, we move towards proving the multiplicity one theorem for $\op{GL}_n$.

\subsection{Whittaker Models}
Let's state our target more precisely.
\begin{theorem}[Multiplicity one]
	Fix a reductive group $G$ over a global field $F$. Then every irreducible admissible representation $\pi$ of $G$ appears at most once in the space of cusp forms of $G$.
\end{theorem}
One has a stronger result for $\op{GL}_n$.
\begin{theorem}[Strong multiplicity one]
	Fix a global field $F$, and choose cuspidal automorphic forms $\pi$ and $\pi'$ of $\op{GL}_n$. If $\pi_v\cong\pi_v'$ for all places $v$ away from a finite set, then $\pi\cong\pi'$.
\end{theorem}
We will need some more local representation theory. In particular, we will want some Whittaker models. These arise from Whittaker functionals.
\begin{notation}
	Fix a split reductive group $G$ over a local nonarchimedean field $F$, and choose a nontrivial character $\psi$ on $F$. Let $\Pi\subseteq\Phi$ be the collection of simple roots induced by a choice of Borel $B\subseteq G$, and then we define $\psi_N$ to be the composite of the projections
	\[N(F)\to\prod_{\alpha\in\Pi}N_\alpha(F)\cong\prod_{\alpha\in\Pi}F\stackrel\prod\to F\stackrel\psi\to\CC^\times.\]
\end{notation}
\begin{remark}
	Here, the map $N\to\prod_{\alpha\in\Pi}N_\alpha$ is the projection onto the simple root spaces. It can be seen to be a homomorphism because it factors through $N/[N,N]$ (because $[N,N]$ does not have anything in the simple root spaces).
\end{remark}
\begin{example}
	For $G=\op{GL}_n$, we have $N$ equal to the subgroup of strictly upper-triangular matrices. Then we can take our character $\psi_N$ to be
	\[\psi_N(u)\coloneqq\psi\Bigg(\sum_{i=1}^{n-1}u_{i,i+1}\Bigg)\]
\end{example}
\begin{definition}[Whittaker functional]
	Fix a split reductive group $G$ over a local nonarchimedean field $F$, and choose a smooth irreducible representation $\pi\colon G(F)\to\op{GL}(V)$. Then a \textit{Whittaker functional} on $(\pi,V)$ is a linear functional $\Lambda$ on $V$ for which
	\[\Lambda(\pi(u)\cdot v)=\psi_N(u)\Lambda(v)\]
	for all $u\in N(F)$ and $v\in V$.
\end{definition}
\begin{remark}
	By definition, $\Lambda$ is an $N(F)$-invariant functional in $V^*=\op{Hom}(V,\CC)$, where $N(F)$ acts on $V$ by $\pi$ and on $\CC$ by the character $\psi_N$. However, $\Lambda$ is not required to be smooth!
\end{remark}
\begin{remark} \label{rem:whittaker-to-embed}
	By Frobenius reciprocity, the data of a Whittaker functional is the data of an element in the space
	\[\op{Hom}_{N(F)}(V,\psi_N)\cong\op{Hom}_{G(F)}\left(V,\op{Ind}_{N(F)}^{G(F)}\psi_N\right),\]
	where $\op{Ind}_{N(F)}^{G(F)}\psi_N$ can be described as smooth functions $\varphi\colon G(F)\to\CC$ satisfying $\varphi(ng)=\psi_N(n)\varphi(g)$ for all $n\in N$ and $g\in G$. (We are using the fact that $N$ and $G$ are unimodular.) In short, one can take $\rho\colon V\to\op{Ind}_{N(F)}^{G(F)}\psi_N$ and produce the map $\Lambda(v)\coloneqq\rho(v)(1)$; it is not hard to define an inverse isomorphism as well.
\end{remark}

\subsection{Existence of Whittaker Models}
By \Cref{rem:whittaker-to-embed}, the existence of a nonzero Whittaker functional amounts to the data of an embedding into $\op{Ind}_{N(F)}^{G(F)}\psi_N$. Thus, it may appear that these functionals are rather restrictive objects.
\begin{theorem} \label{thm:whittaker-exists}
	Fix a split reductive group $G$ over a local nonarchimedean field $F$, and choose a smooth irreducible representation $\pi\colon G(F)\to\op{GL}(V)$.
	\begin{listalph}
		\item The space of Whittaker functionals is at most one-dimensional.
		\item The representation $\pi$ usually admits a nonzero Whittaker functional.
	\end{listalph} 
\end{theorem}
\begin{remark}
	By \Cref{rem:whittaker-to-embed}, we are equivalently asserting that each smooth irreducible representation appears at most once in $\op{Ind}_{N(F)}^{G(F)}\psi_N$, and this is possible for most smooth irreducible $\pi$. One could imagine checking the first property using something like Mackey theory.
\end{remark}
\begin{example}
	Let's make the statement of (b) precise when $G=\op{GL}_2$. We have two cases.
	\begin{itemize}
		\item It turns out that finite-dimensional smooth representations factor through $\det\colon\mathrm{GL}_2\to\mathbb G_m$, so the irreducible ones are one-dimensional. In this case, $N$ vanishes on these representations, so there are no Whittaker functionals.
		\item However, all infinite-dimensional smooth representations admit Whittaker functionals. We will show a special case in \Cref{exe:whittaker-for-unr-induction} below.
	\end{itemize}
\end{example}
\begin{example}
	Roughly speaking, the ``largest'' smooth representations of $\op{GL}_n(F)$ admit Whittaker functionals, so \Cref{thm:whittaker-exists} basically asserts that all these largest smooth representations are typical. For example, the parabolic inductions $\op{Ind}_B^G\chi$ are of the largest possible size, and their largest subquotients admit Whittaker functionals. However, for larger parabolic subgroups $P$, the parabolic inductions $\op{Ind}_P^G\chi$ do not admit Whittaker functionals.
\end{example}
\begin{exe} \label{exe:whittaker-for-unr-induction}
	Fix $G\coloneqq\op{GL}_2$, and suppose that $\psi\colon F\to\CC^\times$ has conductor $\OO_F$. Choose distinct unramified characters $\chi_1,\chi_2\colon F^\times\to\CC^\times$, and let $\pi$ be the parabolic induction $V\coloneqq\op{Ind}_B^G\chi$. We show that $V$ admits a one-dimensional space of Whittaker functionals.
\end{exe}
\begin{proof}
	Then $V$ consists of those smooth functions $f\colon G(F)\to\CC$ for which
	\[f\left(\begin{bmatrix}
		a & 0 \\ c & d
	\end{bmatrix}\right)=\left|\frac ad\right|^{1/2}\chi_1(a)\chi_2(a)f(g).\]
	By \Cref{rem:whittaker-to-embed}, we are looking for an embedding $\op{Ind}_B^G(\chi_1,\chi_2)\to\op{Ind}_N^G\psi_N$.
	
	Define $f_0\in V^K$ by $f_0\left(\begin{bsmallmatrix}
		a & 0 \\ c & d
	\end{bsmallmatrix}k\right)\coloneqq\left|a/d\right|^{1/2}\chi_1(a)\chi_2(d)$; it is not hard to check that $f_0$ is well-defined, and we showed on the homework that it is the unique vector in $V^K$ up to scalar. Because $f_0$ is nonzero, it must generate our irreducible representation, so it is enough to describe the image of $f_0$. Well, $f_0$ must map to a $K$-invariant function in $\op{Ind}_N^G\psi_N$, which is the space of functions $\varphi$ on $G(F)$ which are left invariant by $K$ and right invariant by $(N(F),\psi_N)$. Any such function can be uniquely described by its behavior on representatives of $N(F)\backslash G(F)/K$, which can be taken to be the diagonal elements $\op{diag}\left(\varpi^a,\varpi^b\right)$ by the Iwasawa decomposition.

	Now, a calculation can show that $\varphi$ must vanish on $\op{diag}\left(\varpi^a,\varpi^b\right)$ when $a<b$, using the fact that the conductor of $\psi$ is $\OO_F$. Indeed, set $g\coloneqq\op{diag}\left(\varpi^a,\varpi^b\right)$, and suppose that $\varphi(g)\ne0$, and we will show that $a\ge b$. Well, for any $u\in N(F)$, if $g^{-1}ug\in K$, then
	\[\psi_N(u)\varphi(g)=\varphi(ug)=\varphi\left(g\cdot g^{-1}ug\right)=\varphi(g),\]
	so $\psi_N(u)=1$ is forced. We conclude that $N(F)\cap gKg^{-1}$ must be contained in $[N,N]$, which for $\op{GL}_2$ is trivial. On the other hand, for $u=\begin{bsmallmatrix}
		1 & x \\ & 1
	\end{bsmallmatrix}$, we find that $g^{-1}ug=\begin{bsmallmatrix}
		1 & \varpi^{b-a}x \\ & 1
	\end{bsmallmatrix}$, so we see that $a\ge b$ is forced for this element to fail to be in $K$ when $x\ne0$.

	Thus, the previous paragraph shows that the evaluation map
	\[\op{Ind}_N^G(\psi_N)^K\to\CC^{\{(a,b):a\ge b\}}\]
	is an isomorphism. This is not yet one-dimensional, so we will need to work a little harder to pin down the image of $f_0$. Well, note that the spherical Hecke algebra $\mc H_K$ acts on $V^K=\CC f_0$ by some character $\xi\colon\mc H_K\to\CC$. Thus, the image of $\rho_0$ in $\op{Ind}_N^G(\psi_N)^K$ must also be a $\xi$-eigenvector. We will compute the action of $\mc H_K$ on $\op{Ind}_N^G(\psi_N)^G$ on the homework, and one finds that the relevant eigenspaces are one-dimensional. Matching up the image of $f_0$ completes the proof.
\end{proof}

\subsection{Finite Fields}
Because it is technically easier (but has the same sort of ideas), we will only sketch \Cref{thm:whittaker-exists} over finite fields. All the same definitions go through.
\begin{definition}
	Fix a split reductive group $G$ over a finite field $\FF_q$, and let $N\subseteq G$ be the unipotent radical of the Borel. Then a \textit{Whittaker functional} on a representation $V$ of $G$ is an element of the space
	\[\op{Hom}_{N(F)}(V,\psi_N).\]
\end{definition}
\begin{remark} \label{rem:whittaker-by-embed}
	By Frobenius reciprocity, the data of a Whittaker functional is equivalent to an element of the space $\op{Hom}_{G(F)}\left(V,\op{Ind}_{N(F)}^{G(F)}\psi_N\right)$.
\end{remark}
\begin{theorem}[Gelfand--Graev]
	Fix a reductive group $G$ over a finite field $\FF_q$. Every irreducible representation of $G(\FF_q)$ admits at most one Whittaker functional up to scalar.
\end{theorem}
\begin{proof}
	We identify all groups by their $\FF_q$-points throughout. By \Cref{rem:whittaker-by-embed}, it is enough to show that the representation $\op{Ind}_N^G\psi_N$ is multiplicity-free. Equivalently, we may show that the ring $\op{End}\left(\op{Ind}_N^G\psi_N\right)$ is commutative. By Frobenius reciprocity, this vector space is given by functions $f\colon G\to\CC$ for which
	\[f(u_1gu_2)=\psi_N(u_1u_2)f(g)\]
	for all $g\in G$ and $u_1,u_2\in N$. A direct calculation shows that the ring operation on this space of functions is given by convolution: one has
	\[(f_1*f_2)(g)\coloneqq\sum_{x\in N\backslash G}f_1\left(gx^{-1}\right)f(x).\]
	We will show that this convolution algebra is commutative via the Gelfand trick. Indeed, we will hunt for an anti-involution $\iota$ on $G$ which preserves the subgroup $N$, preserves the character $\psi_N$, and fixes representatives for the double cosets $N\backslash G/N$. Then the map $f\mapsto f^\iota$, where $f^\iota(g)\coloneqq f(\iota g)$, fixes a basis of $N\backslash G/N$ while switching the multiplication, thereby forcing commutativity.

	To this end, we will use
	\[\iota(g)\coloneqq w_0\tau(g)w_0,\]
	where $w_0$ is the longest element of the Weyl group. By construction, $w_0$ flips $N$ to its opposite (it negates all the roots), so $\iota$ preserves $N$, and we can see that it also preserves $\psi_N$.

	However, $\iota$ does not fix representatives of our double cosets of $N\backslash G/N$: indeed, elements of $T$. Luckily, we only need this to be true on representatives for a basis of $\op{Ind}_N^G\psi_N$. For example, for some $t\in T$, we see that
	\[\psi_N(u)f(t)=f(t)\psi_N\left(t^{-1}ut\right).\]
	Thus, having $f(t)\ne0$ forces $\psi_N(u)=\psi_N\left(t^{-1}ut\right)$ for all $u$. The left-hand side is $\psi$ of $\sum_iu_{i,i+1}$, but the right-hand side is $\psi$ of $\sum_ia_i^{-1}a_{i+1}u_{i,i+1}$, so this equality is only possible when $t$ is a scalar.

	To make this work in general, we use the Bruhat decomposition, which tells us that
	\[N\backslash G/N=\bigsqcup_{w\in W}\bigcup_{t\in T}NwtN.\]
	Above, we showed that the piece with $w=1$ has functions only with support when $t$ is a scalar. One can do a similar analysis for other Weyl elements. For example, when $G=\op{GL}_2$, it only remains to worry about $w_0T$, which is already preserved by $\iota$ anyway. For $G=\op{GL}_n$ in general, one finds that Weyl elements (up to permutation) look like $w_Jw_0$, where $w_J$ is some direct sum of smaller antidiagonal Weyl elements; then $t$ must be a similar direct sum of scalar matrices, which is then still fixed by $\iota$.
	
	In general, we have to study the action of $T$ on the various root spaces of $N$ normalized by $w\in W$. Here is the idea: we are hunting for those double cosets $NwtN$ for which there is a function $f$ supported on this double coset. This means that $\psi_N(u)=1$ whenever $t^{-1}w^{-1}uwt\in N$. This condition basically means that $u$ is supported in the subgroup generated by the root spaces $N_\alpha$ for which $w$ sends $\alpha$ to another positive root. Then having $f(tw)\ne0$ means that the adjoint action by $t$ on these root spaces needs to be trivial. I suspect that one can work out that this implies that $\iota$ fixes the double coset $NwtN$.
\end{proof}
\begin{remark}
	Over $p$-adic fields, the same idea works, but we have to work with distributions instead of functions.
\end{remark}

\end{document}